\documentclass{article}

\usepackage{tabularx}
\usepackage{booktabs}
\usepackage{longtable}
\usepackage[numbers,square,sort&compress]{natbib}
\usepackage[letterpaper,left=0.8in,top=1.0in,right=1.0in,bottom=1.0in]{geometry}

\title{Reflection and Traceability Report on \progname}

\author{\authname}

\date{}

\input{../Comments.text}
\input{../Common.text}

\begin{document}

\maketitle

\newpage

\tableofcontents

\listoftables %if appropriate

\newpage

% \plt{Reflection is an important component of getting the full benefits from a
% learning experience.  Besides the intrinsic benefits of reflection, this
% document will be used to help the TAs grade how well your team responded to
% feedback.  Therefore, traceability between Revision 0 and Revision 1 is and
% important part of the reflection exercise.  In addition, several CEAB (Canadian
% Engineering Accreditation Board) Learning Outcomes (LOs) will be assessed based
% on your reflections.}

\section{Changes in Response to Feedback}

% \plt{Summarize the changes made over the course of the project in response to
% feedback from TAs, the instructor, teammates, other teams, the project
% supervisor (if present), and from user testers.}

% \plt{For those teams with an external supervisor, please highlight how the feedback 
% from the supervisor shaped your project.  In particular, you should highlight the 
% supervisor's response to your Rev 0 demonstration to them.}

% \plt{Version control can make the summary relatively easy, if you used issues
% and meaningful commits.  If your feedback is in an issue, and you responded in
% the issue tracker, you can point to the issue as part of explaining your
% changes.  If addressing the issue required changes to code or documentation, you
% can point to the specific commit that made the changes.  Although the links are
% helpful for the details, you should include a label for each item of feedback so
% that the reader has an idea of what each item is about without the need to click
% on everything to find out.}

% \plt{If you were not organized with your commits, traceability between feedback
% and commits will not be feasible to capture after the fact.  You will instead
% need to spend time writing down a summary of the changes made in response to
% each item of feedback.}

% \plt{You should address EVERY item of feedback.  A table or itemized list is
% recommended.  You should record every item of feedback, along with the source of
% that feedback and the change you made in response to that feedback.  The
% response can be a change to your documentation, code, or development process.
% The response can also be the reason why no changes were made in response to the
% feedback.  To make this information manageable, you will record the feedback and
% response separately for each deliverable in the sections that follow.}

% \plt{If the feedback is general or incomplete, the TA (or instructor) will not
% be able to grade your response to feedback.  In that case your grade on this
% document, and likely the Revision 1 versions of the other documents will be
% low.} 

\subsection{Problem Statement, Goals, and Development Plan} %------------------------------------------------------------------------

Table~\ref{Tbl-ProbState} summarizes the feedback received from our peers, the TA, the instructor and other sources.  The table includes links to the related issues and pull requests, along with a summary of our response/changes.

\begin{longtable}{|p{2cm}|p{4cm}|p{5cm}|p{4cm}|}
\caption{Feedback Responses and Traceability for Problem Statement, Goals, Development Plan}
\label{Tbl-ProbState}\\

\hline
\textbf{Feedback Source} & \textbf{Summary} & \textbf{Changes Made} & \textbf{Related Issues} \\ 
\hline
\endhead

Peer & Learning outcome extraction lacks a measurable success criteria & Added measurable success criteria to the learning outcome extraction & \href{https://github.com/shivyam/cas-course-optimization/issues/32}{IS\#32}, \href{https://github.com/shivyam/cas-course-optimization/commit/2a966de}{Commit\#2a966de} \\
\hline
Peer & Word document support is not explicitly addressed in the schedule & Clarified that initially only PDF support will be implemented and Word document support will be documented as future work. & \href{https://github.com/shivyam/cas-course-optimization/issues/33}{IS\#33}, \href{https://github.com/shivyam/cas-course-optimization/commit/f07d482}{Commit\#f07d482} \\
\hline
TA & Need more details in the team charter: How is "repeatedly misses a commitment" defined? What's the exact amount? What would work balancing look like? & Updated team charter to include examples and more details. & \href{https://github.com/shivyam/cas-course-optimization/issues/94}{IS\#94}, \href{https://github.com/shivyam/cas-course-optimization/pull/117}{PR\#117} \\
\hline
Peer & The PoC fallback could leave the CP-SAT constraint risk unresolved & Prioritized a reduced CP-SAT model and clarified that any scoring fallback must document the unresolved constraints and evaluate them before full implementation. & \href{https://github.com/shivyam/cas-course-optimization/issues/35}{IS\#35} \\
\hline
Peer & CEAB validation is not clearly demonstrated in the PoC & Added one representative CEAB graduate-attribute check to the PoC and clarified that broader CEAB coverage will be evaluated afterward. & \href{https://github.com/shivyam/cas-course-optimization/issues/34}{IS\#34} \\
\hline
TA & The PoC plan is too ambitious and does not clearly prioritize its main risks & Made reliable data extraction and alternative scheduling dates the primary PoC focus, while limiting policy, CEAB, explanation, and other features to small demonstrations or follow-up work. & \href{https://github.com/shivyam/cas-course-optimization/issues/93}{IS\#93} \\
\hline
Peer & The PoC integrates too many components (storage, calendar, conflict
detection, alternative dates, LLM explanations, compliance checks) to evaluate
the main risks before the demonstration & Narrowed the PoC in the Development
Plan~\citep{TeamBVSH2026DP} to one outline going through the extraction
pipeline: PDF upload, PyMuPDF text extraction, a local LLM returning
schema-checked JSON, and a frontend table showing each value with its source
passage. Listed the calendar, conflict detection, CP-SAT suggestions, policy and
CEAB checks, LLM explanations, login and export as not part of the
demonstration, and gave a fallback for each risk. This replaces the PoC scope
described for
\href{https://github.com/shivyam/cas-course-optimization/issues/34}{IS\#34},
\href{https://github.com/shivyam/cas-course-optimization/issues/35}{IS\#35} and
\href{https://github.com/shivyam/cas-course-optimization/issues/93}{IS\#93},
since date suggestions became an extra after the meeting with the department. &
\href{https://github.com/shivyam/cas-course-optimization/issues/36}{IS\#36}, \href{https://github.com/shivyam/cas-course-optimization/pull/140}{PR\#140} \\
\hline
TA & Expected Technology needs clearer justification beyond ``X tool/technology
does Y'' (Eval Tools, LO13) & Rewrote the Expected Technology section of the
Development Plan~\citep{TeamBVSH2026DP} so that each tool states how it is used,
why it was chosen over the alternatives considered (for example, PyMuPDF over
pdfplumber based on the model-selection spike results, Ollama because the model
must run on CAS servers at no cost), and its limitations. Added the tools that
were missing (PyMuPDF, Ollama, the candidate open-weight models, OR-Tools
CP-SAT), and marked vLLM, CUDA, Transformers and PyTorch as conditional. & TA
rubric feedback, \href{https://github.com/shivyam/cas-course-optimization/pull/140}{PR\#140} \\
\hline
TA & Keep the quantified success criteria consistent with later documents such
as the SRS & Raised the targets in Goals 2 and 5 of the Problem
Statement~\citep{TeamBVSH2026PS} from 90\% to 95\% to match NFR1 and Appendix C
of the SRS~\citep{TeamBVSH2026SRS}, added a 95\% target for learning outcomes
(with their graduate attributes) to both documents, and changed ``Software
Engineering'' to all CAS undergraduate programs in Goals 1 and 3 to match the
SRS cohorts. &
\href{https://github.com/shivyam/cas-course-optimization/issues/91}{IS\#91}, \href{https://github.com/shivyam/cas-course-optimization/pull/140}{PR\#140} \\
\hline
TA & The predefined outcomes and datasets lack validation; define the ground
truth & Added a ``Test dataset and ground truth'' paragraph to the Goals: real
and mock outlines, and an answer key built by the team that lists every expected
value, missing field, conflict and compliance issue, double-checked against the
outlines before it is used for scoring. &
\href{https://github.com/shivyam/cas-course-optimization/issues/91}{IS\#91}, \href{https://github.com/shivyam/cas-course-optimization/pull/140}{PR\#140} \\
\hline
TA & Has thought been given to courses outside the department (e.g.\ math,
statistics, engineering) that still affect the schedule? & Added a ``Courses
covered'' paragraph to the Goals: required courses from other departments and
first-year engineering courses are not modelled because CAS does not receive or
control their outlines, so the workload shown is a lower bound. Added the
matching assumption A8 to the SRS~\citep{TeamBVSH2026SRS}. &
\href{https://github.com/shivyam/cas-course-optimization/issues/91}{IS\#91}, \href{https://github.com/shivyam/cas-course-optimization/pull/140}{PR\#140} \\
% \hline
% Other & summary of feedback & summary of changes made & \href{link}{IS\#}, \href{link}{PR\#}, ... \\
\hline
\end{longtable}

\subsection{SRS and Hazard Analysis}

The SRS and Hazard Analysis deliverables are closely related, so we will
summarize the feedback received for these two deliverables together in the
sections that follow.

\subsubsection{SRS}

Table~\ref{Tbl-SRS} summarizes the feedback received from our peers, the TA, the
instructor and other sources.  The table includes links to the related issues
and pull requests, along with a summary of our response/changes.

\begin{longtable}{|p{2cm}|p{4cm}|p{5cm}|p{4cm}|}
\caption{Feedback Responses and Traceability for SRS}
\label{Tbl-SRS}\\

\hline
\textbf{Feedback Source} & \textbf{Summary} & \textbf{Changes Made} & \textbf{Related Issues} \\ 
\hline
\endhead

Peer & Session timeout protection is not formally specified in the SRS and Hazard Analysis & Added a requirement specifying that authenticated sessions should expire after a defined period of inactivity in the SRS and updated the Hazard Analysis accordingly & \href{https://github.com/shivyam/cas-course-optimization/issues/86}{IS\#86}, \href{https://github.com/shivyam/cas-course-optimization/pull/118}{PR\#118} \\
\hline
Peer & Model changes could reduce extraction accuracy without a formal regression requirement & Added a requirement to re-run the representative accuracy dataset and confirm the thresholds before deploying a new or updated model. & \href{https://github.com/shivyam/cas-course-optimization/issues/84}{IS\#84} \\
\hline
Peer & CSV export did not specify protection against spreadsheet formula interpretation & Updated the CSV export requirement to escape values beginning with formula characters. & \href{https://github.com/shivyam/cas-course-optimization/issues/85}{IS\#85} \\
\hline
Peer & The SRS does not say how model outputs that fail validation or contain unexpected instructions are rejected (prompt injection) & Added AIR7 (output validation) to the SRS~\citep{TeamBVSH2026SRS}: outline text is given to the model only as data and output is never acted on as an instruction; malformed output is rejected and the outline flagged for manual entry; values without a matching source passage or out of range are treated as missing; policy violations are not rejected; a model is only used if its validated output is unchanged on the prompt-injection test set. Also added a prompt-injection row to Known Limitations, a definition of prompt injection, a rationale item, and AIR7 to the traceability tables. & \href{https://github.com/shivyam/cas-course-optimization/issues/82}{IS\#82}, \href{https://github.com/shivyam/cas-course-optimization/pull/140}{PR\#140} \\
\hline
TA & Keep the success criteria consistent between the Problem Statement and the SRS; consider courses outside the department & Added the 95\% learning outcome accuracy target to the SRS~\citep{TeamBVSH2026SRS} (NFR1, Evaluation Metrics, \textit{LO\_ACCURACY} in Appendix C), changed the source of the accuracy constants to the goals in the Problem Statement~\citep{TeamBVSH2026PS}, updated the accuracy rationale, stated that the answer key is double-checked, and added assumption A8 (only CAS courses are modelled). & \href{https://github.com/shivyam/cas-course-optimization/issues/91}{IS\#91}, \href{https://github.com/shivyam/cas-course-optimization/pull/140}{PR\#140} \\
\hline
% TA & summary of feedback & summary of changes made & \href{link}{IS\#}, \href{link}{PR\#}, ... \\
% \hline
% Instructor & summary of feedback & summary of changes made & \href{link}{IS\#}, \href{link}{PR\#}, ... \\
% \hline
% Other & summary of feedback & summary of changes made & \href{link}{IS\#}, \href{link}{PR\#}, ... \\
% \hline
\end{longtable}

\subsubsection{Hazard Analysis}

Table~\ref{Tbl-Hazard} summarizes the feedback received from our peers, the TA,
the instructor and other sources.  The table includes links to the related
issues and pull requests, along with a summary of our response/changes.

\begin{longtable}{|p{2cm}|p{4cm}|p{5cm}|p{4cm}|}
\caption{Feedback Responses and Traceability for Hazard Analysis}
\label{Tbl-Hazard}\\

\hline
\textbf{Feedback Source} & \textbf{Summary} & \textbf{Changes Made} & \textbf{Related Issues} \\ 
\hline
\endhead

Peer & Backup and recovery safeguards are not clearly defined & Clarified H2-1 by assigning backups and recovery to CAS IT, requiring the project team to document application data and verify representative restores, and adding a backup and restore requirement to the roadmap. & \href{https://github.com/shivyam/cas-course-optimization/issues/83}{IS\#83}, \href{https://github.com/shivyam/cas-course-optimization/pull/123}{PR\#123}, ... \\
\hline
Peer & CSV formula injection prevention was not included in the safety requirements & Added SSR\_CSVExport to formalize escaping dangerous spreadsheet formula prefixes and testing the resulting export. & \href{https://github.com/shivyam/cas-course-optimization/issues/85}{IS\#85} \\
\hline
Peer & Model changes lacked a defined regression testing requirement & Added SSR\_ModelChangeTest to require accuracy testing and recording the model version before deployment. & \href{https://github.com/shivyam/cas-course-optimization/issues/84}{IS\#84} \\
\hline
Peer & The output validation actions for prompt injection (H4-5) are marked ``New'' but are not in the Safety and Security Requirements & Added SSR13 (SSR\_OutputValidation) to Table 5 of the Hazard Analysis~\citep{TeamBVSH2026HA} and linked it to AIR7 in the SRS~\citep{TeamBVSH2026SRS}; replaced ``New: output validation'' in row H4-5 with SSR13 and AIR7; added AIR7 to row H4-1; and listed AIR7 among the SRS safeguards in Section 6. This also responds, for H4-5, to the TA comment that ``New'' requirements should be added to Table 5. & \href{https://github.com/shivyam/cas-course-optimization/issues/82}{IS\#82}, \href{https://github.com/shivyam/cas-course-optimization/pull/140}{PR\#140} \\
% \hline
% TA & summary of feedback & summary of changes made & \href{link}{IS\#},
% \href{link}{PR\#}, ... \\
% \hline
% Instructor & summary of feedback & summary of changes made & \href{link}{IS\#},
% \href{link}{PR\#}, ... \\
% \hline
% Other & summary of feedback & summary of changes made & \href{link}{IS\#},
% \href{link}{PR\#}, ... \\
\hline
\end{longtable}

\subsection{Design and Design Documentation}

The design and design documentation deliverables include the software
architecture and detailed design.  We will summarize the feedback received for
each of these deliverables separately in the sections that follow.

\subsubsection{Software Architecture}

Table~\ref{Tbl-Architecture} summarizes the feedback received from our peers,
the TA, the instructor and other sources.  The table includes links to the
related issues and pull requests, along with a summary of our response/changes.

\begin{longtable}{|p{2cm}|p{4cm}|p{5cm}|p{4cm}|}
\caption{Feedback Responses and Traceability for Software Architecture}
\label{Tbl-Architecture}\\

\hline
\textbf{Feedback Source} & \textbf{Summary} & \textbf{Changes Made} & \textbf{Related Issues} \\ 
\hline
\endhead

Peer & summary of feedback & summary of changes made & \href{link}{IS\#}, \href{link}{PR\#}, ... \\
\hline
TA & summary of feedback & summary of changes made & \href{link}{IS\#}, \href{link}{PR\#}, ... \\
\hline
Instructor & summary of feedback & summary of changes made & \href{link}{IS\#}, \href{link}{PR\#}, ... \\
\hline
Other & summary of feedback & summary of changes made & \href{link}{IS\#}, \href{link}{PR\#}, ... \\
\hline
\end{longtable}

\subsubsection{Detailed Design}

Table~\ref{Tbl-DetailedDesign} summarizes the feedback received from our peers,
the TA, the instructor and other sources.  The table includes links to the
related issues and pull requests, along with a summary of our response/changes.

\begin{longtable}{|p{2cm}|p{4cm}|p{5cm}|p{4cm}|}
\caption{Feedback Responses and Traceability for Detailed Design}
\label{Tbl-DetailedDesign}\\

\hline
\textbf{Feedback Source} & \textbf{Summary} & \textbf{Changes Made} & \textbf{Related Issues} \\ 
\hline
\endhead

Peer & summary of feedback & summary of changes made & \href{link}{IS\#}, \href{link}{PR\#}, ... \\
\hline
TA & summary of feedback & summary of changes made & \href{link}{IS\#}, \href{link}{PR\#}, ... \\
\hline
Instructor & summary of feedback & summary of changes made & \href{link}{IS\#}, \href{link}{PR\#}, ... \\
\hline
Other & summary of feedback & summary of changes made & \href{link}{IS\#}, \href{link}{PR\#}, ... \\
\hline
\end{longtable}

\subsection{VnV Plan and Report}

The VnV Plan and Report deliverables are closely related, so we will summarize
the feedback received for these two deliverables together in the sections that
follow.

\subsubsection{VnV Plan}
Table~\ref{Tbl-VnVPlan} summarizes the feedback received from our peers, the TA, the instructor and other sources.  The table includes links to the related issues and pull requests, along with a summary of our response/changes.

\begin{longtable}{|p{2cm}|p{4cm}|p{5cm}|p{4cm}|}
\caption{Feedback Responses and Traceability for VnV Plan}
\label{Tbl-VnVPlan}\\

\hline
\textbf{Feedback Source} & \textbf{Summary} & \textbf{Changes Made} & \textbf{Related Issues} \\ 
\hline
\endhead

Peer & A new output validation requirement (prompt injection) needs to be verified & Added system test ST\_OutputValidation to the VnV Plan~\citep{TeamBVSH2026VnVP} (injected instructions, malformed output, values without a source passage or out of range) and traced AIR7 and SSR13 to it in the traceability table. & \href{https://github.com/shivyam/cas-course-optimization/issues/82}{IS\#82}, \href{https://github.com/shivyam/cas-course-optimization/pull/140}{PR\#140} \\
\hline
Peer & summary of feedback & summary of changes made & \href{link}{IS\#}, \href{link}{PR\#}, ... \\
\hline
TA & summary of feedback & summary of changes made & \href{link}{IS\#}, \href{link}{PR\#}, ... \\
\hline
Instructor & summary of feedback & summary of changes made & \href{link}{IS\#}, \href{link}{PR\#}, ... \\
\hline
Other & summary of feedback & summary of changes made & \href{link}{IS\#}, \href{link}{PR\#}, ... \\
\hline
\end{longtable}

\subsubsection{VnV Report}

Table~\ref{Tbl-VnVReport} summarizes the feedback received from our peers, the
TA, the instructor and other sources.  The table includes links to the related
issues and pull requests, along with a summary of our response/changes.

\begin{longtable}{|p{2cm}|p{4cm}|p{5cm}|p{4cm}|}
\caption{Feedback Responses and Traceability for VnV Report}
\label{Tbl-VnVReport}\\

\hline
\textbf{Feedback Source} & \textbf{Summary} & \textbf{Changes Made} & \textbf{Related Issues} \\ 
\hline
\endhead

Peer & summary of feedback & summary of changes made & \href{link}{IS\#}, \href{link}{PR\#}, ... \\
\hline
TA & summary of feedback & summary of changes made & \href{link}{IS\#}, \href{link}{PR\#}, ... \\
\hline
Instructor & summary of feedback & summary of changes made & \href{link}{IS\#}, \href{link}{PR\#}, ... \\
\hline
Other & summary of feedback & summary of changes made & \href{link}{IS\#}, \href{link}{PR\#}, ... \\
\hline
\end{longtable}

\section{Custom Documents (CDs)}

\plt{Summarize the CDs that were tackled by this project.  CDs can include
usability testing, code walkthroughs, user documentation, formal proof,
GenderMag personas, Design Thinking, etc.  CDs should have already been approved
by the course instructor as included in your problem statement.  The CDs should
match what is documented in TeamComposition.csv.}

\plt{Provide links to the documentation of the CDs}

\section{Design Iteration (LO10 (Prototypes))}

\plt{Explain how you arrived at your final design and implementation.  How did
the design evolve from the first version to the final version?} 

\plt{Don't just say what you changed, say why you changed it.  The needs of the
client should be part of the explanation.  For example, if you made changes in
response to usability testing, explain what the testing found and what changes
it led to.}

\section{Design Decisions (LO11) (Justifies)}

\plt{Reflect and justify your design decisions.  How did limitations,
 assumptions, and constraints influence your decisions?  Discuss each of these
 separately.}

\section{Economic Considerations (LO19) (Economic)}

\plt{Is there a market for your product? What would be involved in marketing your 
product? What is your estimate of the cost to produce a version that you could 
sell?  What would you charge for your product?  How many units would you have to 
sell to make money? If your product isn't something that would be sold, like an 
open source project, how would you go about attracting users?  How many potential 
users currently exist?}

\section{Reflection on Project Management (LO20) (Manages)}

\plt{This question focuses on processes and tools used for project management.}

\subsection{How Does Your Project Management Compare to Your Development Plan}

\plt{Did you follow your Development plan, with respect to the team meeting plan, 
team communication plan, team member roles and workflow plan.  Did you use the 
technology you planned on using?}

\subsection{What Went Well?}

\plt{What went well for your project management in terms of processes and 
technology?}

\subsection{What Went Wrong?}

\plt{What went wrong in terms of processes and technology?}

\subsection{What Would you Do Differently Next Time?}

\plt{What will you do differently for your next project?}

\section{Ethics and Equity (LO18) (Equity)}

\plt{Describe how your team applied the principle of equity and universal design
to ensure equitable treatment of all stakeholders.}

\section{Reflection on Capstone}

\plt{This question focuses on what you learned during the course of the capstone project.}

\subsection{Which Courses Were Relevant}

\plt{Which of the courses you have taken were relevant for the capstone project?}

\subsection{Knowledge/Skills Outside of Courses}

\plt{What skills/knowledge did you need to acquire for your capstone project
that was outside the courses you took?}

\newpage

\bibliographystyle{../SRS-IEEE/IEEEtran}
\bibliography{../../refs/References}

\end{document}